
\vspace{0.5cm}
\noindent Identificando as dimensões das variáveis envolvidas:
\begin{align*}
[T_{R}] &= T \\
[m_{1}] &= [m_{2}] = M \\
[r] &= L \\
[G] &= \frac{L^{3}}{M \cdot T^{2}}
\end{align*}

\noindent A equação funcional inicial é dada por:
\[ T_{R} = T_{R}(m_{1}, m_{2}, r, G) \]

\noindent Substituindo as dimensões na equação para encontrar os expoentes:
\[ [T_{R}] = M^{2\alpha} \cdot L^{\beta} \cdot \left(\frac{L^{3}}{M \cdot T^{2}}\right)^{\delta} \]

\noindent Igualando os expoentes das dimensões correspondentes de ambos os lados, montamos o sistema:
\[
\begin{cases}
\text{Para M: } 2\alpha - \delta = 0 \\
\text{Para L: } \beta + 3\delta = 0 \\
\text{Para T: } -2\delta = 1
\end{cases}
\]

\noindent Resolvendo o sistema, obtemos os seguintes valores:
\[ \alpha = -\frac{1}{4}, \quad \beta = \frac{3}{2}, \quad \delta = -\frac{1}{2} \]

\noindent Substituindo esses valores de volta e desenvolvendo a expressão:
\begin{align*}
T_{R} &= M^{2\left(-\frac{1}{4}\right)} \cdot L^{\frac{3}{2}} \cdot \left(\frac{L^{3}}{M \cdot T^{2}}\right)^{\left(-\frac{1}{2}\right)} \\
T_{R} &= M^{\left(-\frac{1}{2}\right)} \cdot L^{\frac{3}{2}} \cdot \left(\frac{M \cdot T^{2}}{L^{3}}\right)^{\left(\frac{1}{2}\right)} \\
T_{R} &= \frac{1}{\sqrt{M}} \cdot \sqrt{L^{3}} \cdot \sqrt{\frac{M \cdot T^{2}}{L^{3}}}
\end{align*}

\noindent Manipulando algebricamente para retornar às variáveis originais:
\begin{align*}
T_{R} &= \sqrt{\frac{r^{3}}{m_{2}}} \cdot T_{R} \cdot \sqrt{\frac{m_{2}}{r^{3}}} \cdot 1 \\
T_{R} &= \sqrt{\frac{r^{3}}{m_{2}}} \cdot T_{R} \cdot \sqrt{\frac{m_{2}}{r^{3}}} \cdot \frac{\sqrt{G}}{\sqrt{G}} \\
T_{R} &= \sqrt{\frac{r^{3}}{G \cdot m_{2}}} \cdot T_{R} \cdot \sqrt{\frac{G \cdot m_{2}}{r^{3}}}
\end{align*}